%% GENERATED-SHARED FRAGMENT -- included by BOTH cv/cv.tex (website CV) and
%% jobmarket/base/cv/cv.tex (extended job-market CV). Edit here once; both update.
%%
%% Paper CONTENT only. Each CV lays these out its own way: the website CV prints one
%% "Working Papers" list with abstracts suppressed; the job-market CV splits the JMP into
%% its own section and prints abstracts via \paperabstract. Layout differs on purpose;
%% titles, statuses, abstracts and presentation lists must not.
%%
%% A status changed here (e.g. a conference acceptance) reaches both CVs at once. That is
%% the whole point: the two files previously disagreed about the synergy paper's status.

% ---------- Job market paper ----------
\newcommand{\jmpTitle}{Accounting-Based Governance in Supplier--Customer Relationships}
\newcommand{\jmpAuthors}{(solo-authored)}
\newcommand{\jmpStatus}{Presenting and gathering feedback.}
\newcommand{\jmpAbstract}{I study the monitoring direction of accounting-based provisions in supplier--customer relationships. Using a large language model (LLM), I identify provisions governing the production, dissemination, and verification of accounting information in supply agreements filed with the SEC. These provisions appear in 70.3\% of the agreements, with customers monitoring suppliers more often than suppliers monitor customers. Supplier-targeted monitoring is more likely under cost-plus pricing and when suppliers have higher financial reporting quality. Customer-targeted monitoring is more likely when customers are financially weaker and when the contract includes royalty-based compensation. After contract initiation, customers that monitor suppliers exhibit lower operating costs per sale, particularly under cost-plus pricing. Suppliers that monitor customers hold less cash after contract initiation. Quasi-experimental evidence shows that supplier-targeted monitoring increases following accounting-quality improvements induced by SOX Section~404, while customer-targeted monitoring increases in response to counterparty-risk shocks from the 2018--2019 Section~301 tariffs. Overall, the findings suggest that the direction of accounting information flows reflects distinct economic incentives on the two sides of the supply relationship.}
\newcommand{\jmpPresented}{16th Financial Reporting Workshop; EAA Annual Conferences; Bocconi University; UC3M International Accounting Symposium; Padua Bolzano Early Researcher Consortium.}

% ---------- Merger synergy disclosures ----------
\newcommand{\synTitle}{Are Merger Synergy Disclosures Credible?}
\newcommand{\synAuthors}{(with \href{https://atifellahie.com/}{Atif Ellahie}, \href{https://www.london.edu/faculty-and-research/faculty-profiles/t/tuna-i}{Irem Tuna} \& \href{https://faculty.unibocconi.eu/robertovincenzi/}{Roberto Vincenzi})}
\newcommand{\synStatus}{Accepted at the 2026 \textit{Review of Accounting Studies} Conference.}
\newcommand{\synAbstract}{We examine the credibility and consequences of synergy disclosures in over 12,000 U.S. merger and acquisition (M\&A) transactions announced between 2004 and 2021. Using press releases and conference call transcripts, we identify both qualitative and quantitative statements about expected synergies. We find that these disclosures are more common in larger deals, when more stock-financing is used, and when the targets are publicly listed. Analyst coverage, institutional ownership, and acquirers’ forecasting experience also help explain variation in disclosure practices. Disclosing firms earn higher announcement returns, particularly when disclosures are more detailed or include numeric estimates. However, these synergy expectations are often not fully realized. More intensive disclosures are associated with lower post-merger operating performance, and numeric estimates predict more frequent and larger goodwill impairments. Our findings suggest that while synergy disclosures are positively received by investors, they tend to be overoptimistic, raising concerns about the credibility of forward-looking statements in the M\&A setting. \href{https://ssrn.com/abstract=6815098}{\textcolor{black}{[Available on SSRN]}}}
\newcommand{\synPresented}{European Accounting Symposium for Young Scholars; EAA Annual Conference; Lancaster University Accounting Research \& Policy Symposium; ESSEC Financial Accounting Symposium; Hawaii Accounting Research Conference; University of Manchester; RCF-ECGI Corporate Finance and Governance Conference; Bocconi University; Review of Accounting Studies Conference (scheduled).}

% ---------- Leverage and operating risk ----------
\newcommand{\levTitle}{Leverage and Operating Risk in Asset Pricing: An Accounting Perspective}
\newcommand{\levAuthors}{(with \href{https://www.lse.ac.uk/accounting/people/peter-pope}{Peter Pope})}
\newcommand{\levStatus}{Presenting and gathering feedback.}
\newcommand{\levAbstract}{Separating operating from financial effects helps capture firm fundamentals, as these two types of risk tend to offset each other in the cross-section. We revisit the Fama-French six-factor model (FF6; 2018) by recomputing factors using operating characteristics while separately accounting for financial leverage. Our six-factor model subsumes the original FF6 model and outperforms the q-factor model (HXZ; 2015) in explaining a large set of documented anomalies. The improvement in capturing operating risk primarily stems from using return on net operating assets and growth in net operating assets as measures of profitability and investment, respectively. The enterprise book-to-market ratio becomes redundant once we consider the other operating characteristics. Additionally, after we control for operating effects, leverage is positively associated with returns, as predicted by theory, suggesting a renewed role for a leverage factor in asset pricing.}
\newcommand{\levPresented}{Lancaster University; EAA Annual Conference; Columbia Business School.}
